\subsection{Inter-Category Complementarities}
\label{subsec:inter-category-complementarities}

%%%%%%%%%%%%%%%%%%%%%%%%%%%%%%%%%%%%%%%%%%%%%%%%%%%%%%%%%%%%%%%%%%%%%%%%%%%%%%%
% 10.6: How INU, KNU, and IDN interact with each other
%%%%%%%%%%%%%%%%%%%%%%%%%%%%%%%%%%%%%%%%%%%%%%%%%%%%%%%%%%%%%%%%%%%%%%%%%%%%%%%

Sections~\ref{subsec:three-categories}--\ref{subsec:reference-points} treated 
INU, KNU, and IDN as separate components. But utility categories do not operate 
in isolation---they \emph{interact}. This section develops the complementarity 
structure between categories.

\subsubsection{The Fundamental Insight: Utility Categories Interact}

\paragraph{Beyond Additive Aggregation.}
A naive model would simply add the three utilities:
\begin{equation}
Q_{naive} = \omega_{INU} \cdot U^{INU} + \omega_{KNU} \cdot U^{KNU} + \omega_{IDN} \cdot U^{IDN}
\end{equation}

But this misses crucial interactions:
\begin{itemize}
\item Does individual success enhance or threaten collective welfare?
\item Does personal gain confirm or violate identity?
\item Does group success strengthen or weaken individual standing?
\end{itemize}

\paragraph{The Complementarity Extension.}
The full model includes interaction terms:

\begin{equation}
\boxed{
Q = \sum_{i} \omega_i \cdot U^i + \sum_{i < j} \gamma_{ij} \cdot U^i \cdot U^j
}
\label{eq:complementarity-model}
\end{equation}

where $\gamma_{ij}$ captures the complementarity between categories $i$ and $j$.

%%%%%%%%%%%%%%%%%%%%%%%%%%%%%%%%%%%%%%%%%%%%%%%%%%%%%%%%%%%%%%%%%%%%%%%%%%%%%%%
\subsubsection{The Three Complementarity Coefficients}
%%%%%%%%%%%%%%%%%%%%%%%%%%%%%%%%%%%%%%%%%%%%%%%%%%%%%%%%%%%%%%%%%%%%%%%%%%%%%%%

\paragraph{The Complementarity Matrix.}
With three categories, there are three pairwise interactions:

\begin{table}[htbp]
\centering
\caption{The Inter-Category Complementarity Matrix}
\label{tab:complementarity-matrix}
\renewcommand{\arraystretch}{1.4}
\begin{tabular}{@{}l|ccc@{}}
\toprule
& INU & KNU & IDN \\
\midrule
INU & --- & $\gamma_{INU,KNU}$ & $\gamma_{INU,IDN}$ \\
KNU & $\gamma_{KNU,INU}$ & --- & $\gamma_{KNU,IDN}$ \\
IDN & $\gamma_{IDN,INU}$ & $\gamma_{IDN,KNU}$ & --- \\
\bottomrule
\end{tabular}
\end{table}

Note: The matrix is symmetric ($\gamma_{ij} = \gamma_{ji}$), so there are 
three unique coefficients.

\paragraph{Interpretation of $\gamma_{ij}$.}

\begin{table}[htbp]
\centering
\caption{Interpreting Complementarity Coefficients}
\label{tab:gamma-interpretation}
\renewcommand{\arraystretch}{1.3}
\begin{tabular}{@{}lll@{}}
\toprule
\textbf{Sign} & \textbf{Name} & \textbf{Meaning} \\
\midrule
$\gamma_{ij} > 0$ & Complement & High $U^i$ \emph{enhances} value of high $U^j$ \\
$\gamma_{ij} < 0$ & Substitute & High $U^i$ \emph{reduces} value of high $U^j$ \\
$\gamma_{ij} = 0$ & Independent & $U^i$ and $U^j$ do not interact \\
\bottomrule
\end{tabular}
\end{table}

%%%%%%%%%%%%%%%%%%%%%%%%%%%%%%%%%%%%%%%%%%%%%%%%%%%%%%%%%%%%%%%%%%%%%%%%%%%%%%%
\subsubsection{$\gamma_{INU,KNU}$: Individual--Collective Complementarity}
%%%%%%%%%%%%%%%%%%%%%%%%%%%%%%%%%%%%%%%%%%%%%%%%%%%%%%%%%%%%%%%%%%%%%%%%%%%%%%%

\paragraph{The Question.}
Does individual gain enhance or diminish collective welfare---and vice versa?

\paragraph{Positive Complementarity ($\gamma_{INU,KNU} > 0$).}
Individual and collective interests align:

\begin{table}[htbp]
\centering
\caption{Positive INU-KNU Complementarity Examples}
\label{tab:inu-knu-positive}
\renewcommand{\arraystretch}{1.3}
\begin{tabular}{@{}lll@{}}
\toprule
\textbf{Situation} & \textbf{Mechanism} & \textbf{Effect} \\
\midrule
Team bonus & My success = team success & Both rise together \\
Cooperative venture & Individual effort benefits group & Mutual reinforcement \\
Family prosperity & My income supports family & Aligned interests \\
Public goods (high trust) & My contribution valued & Individual and collective gain \\
\bottomrule
\end{tabular}
\end{table}

When $\gamma_{INU,KNU} > 0$:
\begin{equation}
\frac{\partial^2 Q}{\partial U^{INU} \partial U^{KNU}} > 0
\end{equation}
Gains in one category make gains in the other \emph{more valuable}.

\paragraph{Negative Complementarity ($\gamma_{INU,KNU} < 0$).}
Individual and collective interests conflict:

\begin{table}[htbp]
\centering
\caption{Negative INU-KNU Complementarity Examples}
\label{tab:inu-knu-negative}
\renewcommand{\arraystretch}{1.3}
\begin{tabular}{@{}lll@{}}
\toprule
\textbf{Situation} & \textbf{Mechanism} & \textbf{Effect} \\
\midrule
Zero-sum competition & My gain = your loss & INU↑ implies KNU↓ \\
Free-riding & I benefit, group pays & Individual gain, collective harm \\
Corruption & Personal enrichment & Individual gain undermines collective \\
Exploitation & Extract from group & Negative-sum interaction \\
\bottomrule
\end{tabular}
\end{table}

When $\gamma_{INU,KNU} < 0$:
\begin{equation}
\frac{\partial^2 Q}{\partial U^{INU} \partial U^{KNU}} < 0
\end{equation}
High individual utility \emph{reduces} the marginal value of collective utility.

\paragraph{Empirical Estimates.}
The sign and magnitude of $\gamma_{INU,KNU}$ depend on context:

\begin{table}[htbp]
\centering
\caption{$\gamma_{INU,KNU}$ by Context}
\label{tab:gamma-inu-knu-context}
\renewcommand{\arraystretch}{1.3}
\begin{tabular}{@{}llc@{}}
\toprule
\textbf{Context} & \textbf{Characteristic} & $\gamma_{INU,KNU}$ \\
\midrule
High-trust cooperative & Aligned incentives & $+0.3$ to $+0.5$ \\
Family/close group & Strong bonds & $+0.2$ to $+0.4$ \\
Competitive market & Zero-sum elements & $-0.1$ to $+0.1$ \\
Corrupt institution & Extraction dominates & $-0.2$ to $-0.4$ \\
Prisoner's dilemma & Structural conflict & $-0.3$ to $-0.5$ \\
\bottomrule
\end{tabular}
\end{table}

%%%%%%%%%%%%%%%%%%%%%%%%%%%%%%%%%%%%%%%%%%%%%%%%%%%%%%%%%%%%%%%%%%%%%%%%%%%%%%%
\subsubsection{$\gamma_{INU,IDN}$: Individual--Identity Complementarity}
%%%%%%%%%%%%%%%%%%%%%%%%%%%%%%%%%%%%%%%%%%%%%%%%%%%%%%%%%%%%%%%%%%%%%%%%%%%%%%%

\paragraph{The Question.}
Does personal success confirm or violate one's identity?

\paragraph{Positive Complementarity ($\gamma_{INU,IDN} > 0$).}
Personal success aligns with identity:

\begin{table}[htbp]
\centering
\caption{Positive INU-IDN Complementarity Examples}
\label{tab:inu-idn-positive}
\renewcommand{\arraystretch}{1.3}
\begin{tabular}{@{}lll@{}}
\toprule
\textbf{Situation} & \textbf{Identity} & \textbf{Effect} \\
\midrule
Entrepreneur succeeds & ``I'm a winner'' & Success confirms identity \\
Artist sells work & ``I'm a real artist'' & Market validates identity \\
Athlete wins & ``I'm a champion'' & Victory confirms self-image \\
Scholar publishes & ``I'm an intellectual'' & Achievement fits narrative \\
\bottomrule
\end{tabular}
\end{table}

When $\gamma_{INU,IDN} > 0$, success feels \emph{more} satisfying because it 
confirms who you are.

\paragraph{Negative Complementarity ($\gamma_{INU,IDN} < 0$).}
Personal success conflicts with identity:

\begin{table}[htbp]
\centering
\caption{Negative INU-IDN Complementarity Examples}
\label{tab:inu-idn-negative}
\renewcommand{\arraystretch}{1.3}
\begin{tabular}{@{}lll@{}}
\toprule
\textbf{Situation} & \textbf{Identity Conflict} & \textbf{Effect} \\
\midrule
``Selling out'' & Artist takes corporate job & INU↑ but IDN↓↓ \\
Betraying class & Working-class person ``goes posh'' & Financial gain, identity loss \\
Coal miner → coder & Manual labor identity & Higher pay, identity violation \\
Inherited wealth & ``Self-made'' identity & Wealth doesn't fit narrative \\
\bottomrule
\end{tabular}
\end{table}

When $\gamma_{INU,IDN} < 0$, success feels \emph{hollow} or even \emph{shameful}
because it violates who you are.

\paragraph{The ``Selling Out'' Phenomenon.}
Negative $\gamma_{INU,IDN}$ explains why people reject lucrative opportunities:
\begin{equation}
\Delta Q = \omega_{INU} \cdot \Delta U^{INU} + \gamma_{INU,IDN} \cdot U^{INU} \cdot \Delta U^{IDN}
\end{equation}

If $\Delta U^{INU} > 0$ (financial gain) but $\Delta U^{IDN} < 0$ (identity violation),
and $|\gamma_{INU,IDN} \cdot \Delta U^{IDN}|$ is large, then $\Delta Q < 0$ despite 
the financial gain.

\paragraph{Empirical Estimates.}

\begin{table}[htbp]
\centering
\caption{$\gamma_{INU,IDN}$ by Situation}
\label{tab:gamma-inu-idn-context}
\renewcommand{\arraystretch}{1.3}
\begin{tabular}{@{}llc@{}}
\toprule
\textbf{Situation} & \textbf{Characteristic} & $\gamma_{INU,IDN}$ \\
\midrule
Success confirms identity & ``This is who I am'' & $+0.3$ to $+0.5$ \\
Neutral (identity not salient) & No identity implications & $0$ to $+0.1$ \\
Success questions identity & ``Is this really me?'' & $-0.1$ to $-0.3$ \\
Success violates identity & ``Selling out'' & $-0.3$ to $-0.6$ \\
\bottomrule
\end{tabular}
\end{table}

%%%%%%%%%%%%%%%%%%%%%%%%%%%%%%%%%%%%%%%%%%%%%%%%%%%%%%%%%%%%%%%%%%%%%%%%%%%%%%%
\subsubsection{$\gamma_{KNU,IDN}$: Collective--Identity Complementarity}
%%%%%%%%%%%%%%%%%%%%%%%%%%%%%%%%%%%%%%%%%%%%%%%%%%%%%%%%%%%%%%%%%%%%%%%%%%%%%%%

\paragraph{The Question.}
Does group success strengthen or threaten personal identity?

\paragraph{Positive Complementarity ($\gamma_{KNU,IDN} > 0$).}
Group and identity reinforce each other:

\begin{table}[htbp]
\centering
\caption{Positive KNU-IDN Complementarity Examples}
\label{tab:knu-idn-positive}
\renewcommand{\arraystretch}{1.3}
\begin{tabular}{@{}lll@{}}
\toprule
\textbf{Situation} & \textbf{Mechanism} & \textbf{Effect} \\
\midrule
``We are miners'' & Group identity & Collective success = identity affirmation \\
National pride & Patriotic identity & Nation's success enhances self-worth \\
Team identification & ``I am a [team] fan'' & Team victory feels personal \\
Religious community & Faith identity & Community flourishing = personal meaning \\
\bottomrule
\end{tabular}
\end{table}

When $\gamma_{KNU,IDN} > 0$, collective success \emph{validates} personal identity.
This is the foundation of group-based identity and nationalism.

\paragraph{Negative Complementarity ($\gamma_{KNU,IDN} < 0$).}
Group demands conflict with personal identity:

\begin{table}[htbp]
\centering
\caption{Negative KNU-IDN Complementarity Examples}
\label{tab:knu-idn-negative}
\renewcommand{\arraystretch}{1.3}
\begin{tabular}{@{}lll@{}}
\toprule
\textbf{Situation} & \textbf{Identity Conflict} & \textbf{Effect} \\
\midrule
Conformity pressure & ``I must fit in'' vs. authenticity & Group success, identity loss \\
Groupthink & Personal values overridden & Collective gain, self-betrayal \\
Cult dynamics & Identity subsumed by group & KNU↑ but IDN↓ \\
Whistleblower dilemma & Group loyalty vs. integrity & Identity demands betraying group \\
\bottomrule
\end{tabular}
\end{table}

When $\gamma_{KNU,IDN} < 0$, serving the group feels like \emph{losing oneself}.

\paragraph{Empirical Estimates.}

\begin{table}[htbp]
\centering
\caption{$\gamma_{KNU,IDN}$ by Context}
\label{tab:gamma-knu-idn-context}
\renewcommand{\arraystretch}{1.3}
\begin{tabular}{@{}llc@{}}
\toprule
\textbf{Context} & \textbf{Characteristic} & $\gamma_{KNU,IDN}$ \\
\midrule
Strong group identity & ``We are one'' & $+0.4$ to $+0.6$ \\
Moderate identification & Member but individual & $+0.1$ to $+0.3$ \\
Individualistic culture & Self > group & $0$ to $-0.1$ \\
Oppressive group & Group demands conformity & $-0.2$ to $-0.4$ \\
\bottomrule
\end{tabular}
\end{table}

%%%%%%%%%%%%%%%%%%%%%%%%%%%%%%%%%%%%%%%%%%%%%%%%%%%%%%%%%%%%%%%%%%%%%%%%%%%%%%%
\subsubsection{The Complementarity Triad: Coherence and Conflict}
%%%%%%%%%%%%%%%%%%%%%%%%%%%%%%%%%%%%%%%%%%%%%%%%%%%%%%%%%%%%%%%%%%%%%%%%%%%%%%%

\paragraph{Utility Coherence.}
The three complementarities together determine whether the utility system 
is coherent or conflicted:

\begin{definition}[Utility Coherence Index]
\label{def:utility-coherence}
\begin{equation}
K^{utility} = \frac{1}{3} \left( \text{sign}(\gamma_{INU,KNU}) + \text{sign}(\gamma_{INU,IDN}) + \text{sign}(\gamma_{KNU,IDN}) \right)
\end{equation}
\end{definition}

\begin{table}[htbp]
\centering
\caption{Utility Coherence States}
\label{tab:utility-coherence-states}
\renewcommand{\arraystretch}{1.3}
\begin{tabular}{@{}cccll@{}}
\toprule
$\gamma_{INU,KNU}$ & $\gamma_{INU,IDN}$ & $\gamma_{KNU,IDN}$ & $K^{utility}$ & \textbf{State} \\
\midrule
+ & + & + & +1 & \textbf{Full coherence} \\
+ & + & $-$ & +1/3 & Partial coherence \\
+ & $-$ & + & +1/3 & Partial coherence \\
$-$ & + & + & +1/3 & Partial coherence \\
+ & $-$ & $-$ & $-$1/3 & Partial conflict \\
$-$ & + & $-$ & $-$1/3 & Partial conflict \\
$-$ & $-$ & + & $-$1/3 & Partial conflict \\
$-$ & $-$ & $-$ & $-$1 & \textbf{Full conflict} \\
\bottomrule
\end{tabular}
\end{table}

\paragraph{Full Coherence (All Positive).}
When all three $\gamma > 0$:
\begin{itemize}
\item Individual success benefits the group
\item Personal achievement confirms identity
\item Group success validates individual identity
\end{itemize}

This is the \emph{ideal state}---all utility categories reinforce each other.
Example: A successful entrepreneur who builds a thriving community business 
that expresses their values.

\paragraph{Full Conflict (All Negative).}
When all three $\gamma < 0$:
\begin{itemize}
\item Individual gain harms the group
\item Personal success violates identity
\item Group demands betray individual values
\end{itemize}

This is a \emph{pathological state}---no choice can satisfy all categories.
Example: A corrupt official who enriches themselves at public expense while 
violating their self-image as honest and serving a system they despise.

%%%%%%%%%%%%%%%%%%%%%%%%%%%%%%%%%%%%%%%%%%%%%%%%%%%%%%%%%%%%%%%%%%%%%%%%%%%%%%%
\subsubsection{Complementarities and Behavioral Change}
%%%%%%%%%%%%%%%%%%%%%%%%%%%%%%%%%%%%%%%%%%%%%%%%%%%%%%%%%%%%%%%%%%%%%%%%%%%%%%%

\paragraph{Why Complementarities Matter for EBF.}
The complementarity structure determines whether behavioral change is 
reinforced or resisted:

\begin{table}[htbp]
\centering
\caption{Complementarities and Change Dynamics}
\label{tab:complementarities-change}
\renewcommand{\arraystretch}{1.3}
\begin{tabular}{@{}lll@{}}
\toprule
\textbf{Complementarity Pattern} & \textbf{Change Dynamic} & \textbf{Implication} \\
\midrule
All positive & Self-reinforcing & Change in one → supports others \\
Mixed & Partial resistance & Some dimensions resist \\
All negative & Self-undermining & Change in one → harms others \\
\bottomrule
\end{tabular}
\end{table}

\paragraph{The Coal Miner Revisited.}
Consider the coal miner's transition to coder through the complementarity lens:

\begin{table}[htbp]
\centering
\caption{Coal Miner Transition: Complementarity Analysis}
\label{tab:miner-complementarity}
\renewcommand{\arraystretch}{1.3}
\begin{tabular}{@{}lll@{}}
\toprule
\textbf{Complementarity} & \textbf{Status Quo (Mining)} & \textbf{After Transition (Coding)} \\
\midrule
$\gamma_{INU,KNU}$ & + (supports community) & $-$ (leaves community) \\
$\gamma_{INU,IDN}$ & + (mining fits identity) & $-$ (coding violates identity) \\
$\gamma_{KNU,IDN}$ & + (``we are miners'') & $-$ (betraying ``us'') \\
\midrule
$K^{utility}$ & +1 (coherent) & $-$1 (conflicted) \\
\bottomrule
\end{tabular}
\end{table}

The transition moves from full coherence to full conflict---even if $U^{INU}$ 
increases, the complementarity collapse makes total welfare decrease.

\paragraph{Intervention Implication.}
Successful behavioral change requires maintaining or rebuilding complementarities:
\begin{itemize}
\item \textbf{Identity bridge:} Frame new behavior as consistent with identity
\item \textbf{Collective transition:} Move the group together, preserving KNU
\item \textbf{Narrative integration:} Help construct a story where change fits
\end{itemize}

%%%%%%%%%%%%%%%%%%%%%%%%%%%%%%%%%%%%%%%%%%%%%%%%%%%%%%%%%%%%%%%%%%%%%%%%%%%%%%%
\subsubsection{Context Effects on Complementarities}
%%%%%%%%%%%%%%%%%%%%%%%%%%%%%%%%%%%%%%%%%%%%%%%%%%%%%%%%%%%%%%%%%%%%%%%%%%%%%%%

\paragraph{Complementarities Are Context-Dependent.}
The $\gamma_{ij}$ coefficients respond to the $\Psi$-dimensions:

\begin{equation}
\gamma_{ij}(\Psi) = \bar{\gamma}_{ij} + \sum_k \Gamma^{\gamma}_{ij,k} \cdot \Psi_k
\end{equation}

\paragraph{Key Context Effects.}

\begin{table}[htbp]
\centering
\caption{Context Effects on Complementarities}
\label{tab:context-complementarities}
\renewcommand{\arraystretch}{1.3}
\begin{tabular}{@{}llll@{}}
\toprule
\textbf{$\Psi$-Dimension} & \textbf{Effect on $\gamma_{INU,KNU}$} & \textbf{Effect on $\gamma_{INU,IDN}$} & \textbf{Effect on $\gamma_{KNU,IDN}$} \\
\midrule
$\Psi_S$ (Trust) & ↑ (alignment increases) & ~ (neutral) & ↑ (group identity strengthens) \\
$\Psi_I$ (Institutions) & ↑ (less zero-sum) & ~ (neutral) & ~ (neutral) \\
$\Psi_M$ (Markets) & ↓ (more competition) & ↓ (selling out easier) & ↓ (individualism rises) \\
$\Psi_T$ (Stability) & ↑ (long-term alignment) & ↑ (identity more stable) & ↑ (group bonds strengthen) \\
\bottomrule
\end{tabular}
\end{table}

\paragraph{Example: Globalization Effect.}
Market integration ($\Psi_M$ ↑) tends to:
\begin{itemize}
\item Decrease $\gamma_{INU,KNU}$: More competition, less collective benefit sharing
\item Decrease $\gamma_{INU,IDN}$: More opportunities to ``sell out''
\item Decrease $\gamma_{KNU,IDN}$: Weaker group identities, more individualism
\end{itemize}

This explains why globalization can increase material welfare ($U^{INU}_F$) 
while decreasing overall wellbeing---the complementarity structure deteriorates.

%%%%%%%%%%%%%%%%%%%%%%%%%%%%%%%%%%%%%%%%%%%%%%%%%%%%%%%%%%%%%%%%%%%%%%%%%%%%%%%
\subsubsection{Summary: Inter-Category Complementarities}
%%%%%%%%%%%%%%%%%%%%%%%%%%%%%%%%%%%%%%%%%%%%%%%%%%%%%%%%%%%%%%%%%%%%%%%%%%%%%%%

\begin{tcolorbox}[colback=gray!5!white, colframe=gray!75!black, title=Section 10.6 Summary]
\textbf{The Extended Welfare Function:}
$$Q = \sum_{i} \omega_i \cdot U^i + \sum_{i < j} \gamma_{ij} \cdot U^i \cdot U^j$$

\textbf{The Three Complementarities:}
\begin{itemize}
\item $\gamma_{INU,KNU}$: Individual--Collective alignment
\item $\gamma_{INU,IDN}$: Individual--Identity alignment
\item $\gamma_{KNU,IDN}$: Collective--Identity alignment
\end{itemize}

\textbf{Interpretation:}
\begin{itemize}
\item $\gamma > 0$: Categories reinforce each other (complement)
\item $\gamma < 0$: Categories undermine each other (substitute)
\item $\gamma = 0$: Categories are independent
\end{itemize}

\textbf{Utility Coherence:}
\begin{itemize}
\item All $\gamma > 0$: Full coherence (ideal state)
\item All $\gamma < 0$: Full conflict (pathological state)
\item Mixed: Partial coherence/conflict
\end{itemize}

\textbf{Behavioral Change Implication:}
\begin{itemize}
\item Successful change requires maintaining complementarity structure
\item Identity bridges, collective transitions, narrative integration
\item Context ($\Psi$) modulates complementarities
\end{itemize}
\end{tcolorbox}

\vspace{1em}
\hrule
\vspace{0.5em}
\noindent\textit{Next section: 10.7 Context Modulation ($\Gamma$-Matrix)}

%%%%%%%%%%%%%%%%%%%%%%%%%%%%%%%%%%%%%%%%%%%%%%%%%%%%%%%%%%%%%%%%%%%%%%%%%%%%%%%
